% English translation of checked_pages/esq_119_020.tex; source: cote 119, site 20, PDF page 21, typescript page 18.
% AI-translated by GPT-6 Astra (OpenAI, ChatGPT Pro mode).
\scanpage{119/20}
\mspage{18}

Faced with this, some are content to shrug their shoulders
with a disillusioned air and bet that nothing can be drawn
from it all except dreams. They forget, or are unaware,
that our science, and every science, would amount to very little
if, since its beginnings, it had not been nourished by the dreams and visions
of those who devote themselves to it with passion.

4.        From the beginning of my reflection on two-dimensional maps,
I was especially interested in the so-called
``\uline{regular}'' maps, that is, those whose automorphism group
acts transitively (and consequently, simply transitively)
on the set of frames. In the oriented case and in terms of
the algebraic-geometric interpretation in the preceding paragraph,
these are the maps that correspond to a \uline{Galois} covering of
the projective line. Very quickly, too, and even before the connection
with algebraic geometry appeared, it also became
necessary not to exclude infinite maps, which in particular enter
naturally as universal coverings of
finite maps. It appears (as an immediate consequence of the ``dictionary''
of maps, extended to maps that need not be finite)
that for every pair of natural numbers $p, q$ $\geq$ 1, there exists up to
(nonunique) isomorphism one and only one 1-\uline{connected} map
of type $(p,q)$, i.e. all of whose vertices have order $p$
and all of whose faces have order $q$, and this map is regular.
It is pinned down by the choice of a frame, and its
automorphism group is then canonically isomorphic to the quotient
of the cartographic group (resp. of the oriented cartographic
group in the oriented case) by the additional relations
\[ \rho_s^{\,p} = \rho_f^{\,q} = 1. \]
The case in which this group is finite is the ``Pythagorean'' case of
regular spherical maps; the case in which it is infinite gives the regular tessellations

\margin{20}

